% Alur fungsi hash: pesan berukuran bebas dipetakan menjadi digest berukuran tetap.
% Dirujuk dari section-01.tex dengan % Alur fungsi hash: pesan berukuran bebas dipetakan menjadi digest berukuran tetap.
% Dirujuk dari section-01.tex dengan % Alur fungsi hash: pesan berukuran bebas dipetakan menjadi digest berukuran tetap.
% Dirujuk dari section-01.tex dengan % Alur fungsi hash: pesan berukuran bebas dipetakan menjadi digest berukuran tetap.
% Dirujuk dari section-01.tex dengan \input{figures/alur-fungsi-hash}.
\begin{figure}[htbp]
    \centering
    \begin{tikzpicture}[
        kotak/.style={draw, rounded corners, align=center, font=\small, inner sep=6pt},
        panah/.style={-{Latex[length=2.4mm]}, thick},
        node distance=1.4cm,
    ]
        \node[kotak, fill=blue!6, minimum width=3.4cm, minimum height=1.5cm] (m)
            {Pesan $M$\\[2pt]\scriptsize panjang bebas};
        \node[kotak, fill=orange!12, right=of m, minimum width=3.2cm, minimum height=1.9cm] (h)
            {Fungsi hash\\[2pt]$h = H(M)$};
        \node[kotak, fill=green!10, right=of h, minimum width=3.2cm, minimum height=1.5cm] (d)
            {Digest $h$\\[2pt]\scriptsize panjang tetap};
        \draw[panah] (m) -- (h);
        \draw[panah] (h) -- (d);
    \end{tikzpicture}
    \caption{Alur kerja fungsi hash: pesan berukuran bebas dipetakan menjadi digest berukuran tetap}
    \label{fig:alur-fungsi-hash}
\end{figure}
.
\begin{figure}[htbp]
    \centering
    \begin{tikzpicture}[
        kotak/.style={draw, rounded corners, align=center, font=\small, inner sep=6pt},
        panah/.style={-{Latex[length=2.4mm]}, thick},
        node distance=1.4cm,
    ]
        \node[kotak, fill=blue!6, minimum width=3.4cm, minimum height=1.5cm] (m)
            {Pesan $M$\\[2pt]\scriptsize panjang bebas};
        \node[kotak, fill=orange!12, right=of m, minimum width=3.2cm, minimum height=1.9cm] (h)
            {Fungsi hash\\[2pt]$h = H(M)$};
        \node[kotak, fill=green!10, right=of h, minimum width=3.2cm, minimum height=1.5cm] (d)
            {Digest $h$\\[2pt]\scriptsize panjang tetap};
        \draw[panah] (m) -- (h);
        \draw[panah] (h) -- (d);
    \end{tikzpicture}
    \caption{Alur kerja fungsi hash: pesan berukuran bebas dipetakan menjadi digest berukuran tetap}
    \label{fig:alur-fungsi-hash}
\end{figure}
.
\begin{figure}[htbp]
    \centering
    \begin{tikzpicture}[
        kotak/.style={draw, rounded corners, align=center, font=\small, inner sep=6pt},
        panah/.style={-{Latex[length=2.4mm]}, thick},
        node distance=1.4cm,
    ]
        \node[kotak, fill=blue!6, minimum width=3.4cm, minimum height=1.5cm] (m)
            {Pesan $M$\\[2pt]\scriptsize panjang bebas};
        \node[kotak, fill=orange!12, right=of m, minimum width=3.2cm, minimum height=1.9cm] (h)
            {Fungsi hash\\[2pt]$h = H(M)$};
        \node[kotak, fill=green!10, right=of h, minimum width=3.2cm, minimum height=1.5cm] (d)
            {Digest $h$\\[2pt]\scriptsize panjang tetap};
        \draw[panah] (m) -- (h);
        \draw[panah] (h) -- (d);
    \end{tikzpicture}
    \caption{Alur kerja fungsi hash: pesan berukuran bebas dipetakan menjadi digest berukuran tetap}
    \label{fig:alur-fungsi-hash}
\end{figure}
.
\begin{figure}[htbp]
    \centering
    \begin{tikzpicture}[
        kotak/.style={draw, rounded corners, align=center, font=\small, inner sep=6pt},
        panah/.style={-{Latex[length=2.4mm]}, thick},
        node distance=1.4cm,
    ]
        \node[kotak, fill=blue!6, minimum width=3.4cm, minimum height=1.5cm] (m)
            {Pesan $M$\\[2pt]\scriptsize panjang bebas};
        \node[kotak, fill=orange!12, right=of m, minimum width=3.2cm, minimum height=1.9cm] (h)
            {Fungsi hash\\[2pt]$h = H(M)$};
        \node[kotak, fill=green!10, right=of h, minimum width=3.2cm, minimum height=1.5cm] (d)
            {Digest $h$\\[2pt]\scriptsize panjang tetap};
        \draw[panah] (m) -- (h);
        \draw[panah] (h) -- (d);
    \end{tikzpicture}
    \caption{Alur kerja fungsi hash: pesan berukuran bebas dipetakan menjadi digest berukuran tetap}
    \label{fig:alur-fungsi-hash}
\end{figure}
